\documentclass{article}
\usepackage{geometry}
\usepackage{graphicx} % Required for inserting images
\usepackage{amsmath, amsthm, amssymb}
\usepackage{parskip}
\newgeometry{vmargin={15mm}, hmargin={24mm,34mm}}
\theoremstyle{definition} 
\newtheorem{definition}{Definition}

\newtheorem{theorem}{Theorem}[section]
\newtheorem{lemma}[theorem]{Lemma}
\newtheorem{proposition}[theorem]{Proposition}
\newtheorem{example}[theorem]{Example}
\newtheorem{corollary}{Corollary}[theorem]
\newcommand{\nil}{\operatorname{nil}}
\newcommand{\ann}{\operatorname{Ann}}

\title{Length of a Module}
\date{April 2025}
\author{Boran Erol}

\begin{document}

\maketitle

% Richard E. Borcherds has a nice proof of the uniqueness of length, I think.
% It's Chapter 3 in Gathmann's commutative algebra notes.

The length of a module is a generalization of the dimension of a vector space.

% Page 155 in ATOCA

\begin{lemma}
    If $ l(M) \leq n $, $ M $ is generated by at most $ n $ elements.
\end{lemma}
\begin{proof}
    We induct on the length. If the length is 0 or 1,
    the result is immediate. Assume the result holds for some
    $ n - 1 $ where $ n \geq 1 $.

    Let $ M $ be a module such that the length of $ M $ is at most $ n $.
    Then, $ M $ has a composition series of length $ n $.
    Then, $ M $ contains a submodule $ N $ of length at most $ n - 1 $
    such that $ M/N $ is simple.

    % TODO: Finish this, but it's easy anyway.
\end{proof}



\end{document}
